%!TEX root = ../../main.tex
\section{리그 오브 레전드와 한타}
\label{sec:intro-game}

리그 오브 레전드(League of Legends, LoL)는 다섯 명씩 두 팀이 소환사의 협곡이라는 대칭 지도에서 겨루어 상대의 넥서스를 파괴하면 이기는 게임이다 \cite{riotapi2024matchv5,lolwiki2025summonersrift,lolwiki2026nexus}. 각 플레이어는 챔피언 하나를 조종하며, 챔피언은 미니언과 몬스터를 처치하고 상대 챔피언을 킬하여 얻는 골드와 경험치로 강해진다 \cite{lolwiki2026champion,lolwiki2026gold,lolwiki2026experience}. 챔피언 킬은 킬러와 어시스트에게 골드를 주고 희생자를 사망 시간 동안 경기에서 제외하며 \cite{lolwiki2026kill,lolwiki2026assist,lolwiki2026death}, 포탑·억제기 같은 구조물과 드래곤·전령·바론 같은 엘리트 몬스터는 이를 가져간 팀에 자원과 지도 통제권을 준다 \cite{lolwiki2026turret,lolwiki2026inhibitor,lolwiki2026dragonpit,lolwiki2026riftherald,lolwiki2026baron}. 넥서스는 두 넥서스 포탑과 억제기 하나 이상이 파괴된 뒤에만 공격할 수 있으므로, 경기는 골드·경험치의 우위가 구조물과 오브젝트의 우위로, 다시 넥서스의 파괴로 이어지는 긴 과정이다.

한타(teamfight)는 여러 챔피언이 한 곳에서 싸우는 국면을 가리키는 플레이어의 용어이며 공식 정의는 없다 \cite{lolwiki2026terminology}. 교전에서 챔피언이 죽고 살아남은 쪽이 구조물이나 엘리트 몬스터를 가져가므로, 경기의 흐름은 대체로 교전에서 바뀐다. 교전의 결과는 킬 수만으로 정해지는 양이 아니다. 킬에서 이긴 쪽이 지도 통제권을 잃을 수 있고, 대등하게 교환한 교전이 그 뒤에 가능해진 오브젝트 때문에 경기를 바꿀 수도 있다. 교전의 결과를 어떤 양으로 정의할 것인가가 이 논문의 첫 번째 문제이다(제 \ref{sec:intro-problem} 절).

Riot Games는 모든 경기의 기록을 Match-V5 API로 공개한다 \cite{riotapi2024matchv5}. 경기의 timeline은 서로 다른 해상도의 두 열로 이루어진다 \cite{riotapi2024timeline}. 참가자 프레임은 약 60초마다 열 명의 위치, 레벨, 골드, 체력을 기록하고, 사건은 밀리초 단위의 시각으로 챔피언 킬, 구조물 파괴, 엘리트 몬스터 처치 등을 기록한다. 챔피언 사이에 오간 피해는 사건으로 기록되지 않으며, 챔피언 대 챔피언 전투 가운데 시각과 지도 위치가 모두 기록되는 사건은 챔피언 킬뿐이다. 따라서 공개 자료에서 교전은 킬로부터 구성해야 하고, 어느 시각의 상태는 마지막 프레임과 그때까지의 사건을 합쳐 재구성해야 한다. 표 \ref{tab:terms}는 본 논문이 쓰는 게임·자료 용어를 모은 것이다. Riot은 규칙을 패치로 바꾸며, 본 논문의 코퍼스는 Match-V5 표기 15.14--15.16(공개 번호 25.14--25.16)의 세 패치에 있고 \cite{lolwiki2026v2514,lolwiki2026v2516}, 이후 패치 26.1에서는 일부 오브젝트의 규칙이 바뀌었다 \cite{riot2026patch261}.

%% table 객체: tab:terms — 게임·자료 용어 (docs/tog_manuscript/sec_background.tex Table tab:terms의 국문 요약; 출처는 각 행의 bib 키)
\begin{table}[!t]
  \centering
  \caption{게임과 자료 용어}
  \label{tab:terms}
  \footnotesize
  \begin{tabularx}{\linewidth}{@{}>{\raggedright\arraybackslash}p{3.0cm}X>{\raggedright\arraybackslash}p{3.4cm}@{}}
    \toprule
    용어 & 뜻 & 출처 \\
    \midrule
    소환사의 협곡 & 다섯 명씩 두 팀이 겨루는 대칭 지도; 세 라인과 정글, 강 & \cite{lolwiki2025summonersrift,riotapi2024matchv5} \\
    챔피언 & 플레이어가 조종하는 캐릭터; 팀당 다섯 & \cite{lolwiki2026champion} \\
    골드·경험치 & 미니언·몬스터·킬에서 얻는 자원; 아이템과 레벨로 전환되어 챔피언을 강하게 함 & \cite{lolwiki2026gold,lolwiki2026experience} \\
    킬·어시스트·사망 & 킬은 킬러에게 골드를 줌; 어시스트는 직전 기여자에게 인정; 사망한 챔피언은 일정 시간 경기에서 제외 & \cite{lolwiki2026kill,lolwiki2026assist,lolwiki2026death} \\
    포탑·억제기·넥서스 & 라인을 지키는 구조물; 넥서스 파괴가 승리 조건 & \cite{lolwiki2026turret,lolwiki2026inhibitor,lolwiki2026nexus} \\
    엘리트 몬스터 & 드래곤, 협곡의 전령, 공허 유충, 아타칸, 바론; 처치한 팀에 보상 & \cite{lolwiki2026dragonpit,lolwiki2026riftherald,lolwiki2026voidgrub,lolwiki2026atakhan,lolwiki2026baron} \\
    시야·와드 & 시야 밖의 적 챔피언은 보이지 않음; 와드는 시야를 제공 & \cite{lolwiki2026sight,lolwiki2025ward} \\
    패치 & Riot이 규칙을 바꾸는 단위; Match-V5는 15.N, 공개 번호는 25.N & \cite{riot2025patch25s11,lolwiki2026v2514} \\
    참가자 프레임 & Match-V5 timeline의 약 60초 간격 상태 기록(열 명의 위치·레벨·골드·체력) & \cite{riotapi2024timeline} \\
    사건 & Match-V5 timeline의 밀리초 시각 기록(킬·구조물·몬스터 등); 챔피언 킬만 시각과 위치를 모두 가짐 & \cite{riotapi2024timeline} \\
    교전 (본 논문) & 한타와 소규모 교전을 아우르는 분석 단위. 킬을 기준으로 시간·공간 범위와 참여 조건을 정해 구성한다 & 제 \ref{ch:engagement} 장 \\
    한타 & 교전 가운데 적게 모인 쪽이 네 명 이상인 것. 플레이어는 인원이 많이 모인 맞붙음을 이렇게 부른다 & \cite{lolwiki2026terminology}, 제 \ref{ch:engagement} 장 \\
    소규모 교전 & 교전 가운데 적게 모인 쪽이 두 명 또는 세 명인 것 & 제 \ref{ch:engagement} 장 \\
    \bottomrule
  \end{tabularx}
\end{table}

